% -*- coding: utf-8 -*-
% Bitácora de uso de inteligencia artificial — Intercambiador de Tubos y Coraza
% Documento compilable (tectonic / pdflatex).
\documentclass[11pt,a4paper]{article}

\usepackage[utf8]{inputenc}
\usepackage[T1]{fontenc}
\usepackage{lmodern}
\usepackage[spanish]{babel}
\usepackage[table]{xcolor}
\usepackage{geometry}
\usepackage{fancyhdr}
\usepackage{titlesec}
\usepackage{booktabs}
\usepackage{longtable}
\usepackage{array}
\usepackage{hyperref}

% --- Paleta morada del proyecto -------------------------------------------
\geometry{margin=2.4cm}
\definecolor{morado}{HTML}{5E2B8C}
\definecolor{acento}{HTML}{7E57C2}
\definecolor{seccion}{HTML}{EDE7F6}
\definecolor{gristexto}{HTML}{3B3648}
\hypersetup{colorlinks=true, linkcolor=morado, urlcolor=acento}

\titleformat{\section}{\Large\bfseries\color{morado}}{\thesection}{0.6em}{}[{\color{acento}\titlerule}]
\titleformat{\subsection}{\large\bfseries\color{acento}}{\thesubsection}{0.5em}{}
\renewcommand{\familydefault}{\sfdefault}

\pagestyle{fancy}
\fancyhf{}
\fancyhead[L]{\small\color{morado}\textbf{Boh\'orquez $\cdot$ M\'arquez $\cdot$ Monta\~no $\cdot$ Mor\'e}}
\fancyhead[R]{\small\color{gristexto}Intercambiador de Tubos y Coraza}
\fancyfoot[C]{\small\color{gristexto}\thepage}

\emergencystretch=3em
\begin{document}

% --- Portada ---------------------------------------------------------------
\begin{center}
{\color{acento}\rule{\linewidth}{2pt}}\\[1.2cm]
{\Huge\bfseries\color{morado} Bitácora de uso de\\[0.4em]
inteligencia artificial}\\[0.6em]
{\Large\color{acento} Herramienta computacional educativa para el diseño de\\
intercambiadores de calor de tubos y coraza}\\[0.5em]
{\large\color{gristexto} Método de Kern $\cdot$ Norma TEMA Clase C}\\[0.8cm]
{\color{acento}\rule{\linewidth}{1pt}}\\[0.6cm]
{\large\color{gristexto} Transferencia de Calor / Diseño de Equipos Térmicos
\quad $\cdot$ \quad Octubre de 2026}\\[0.8cm]
{\large\color{gristexto}
\textbf{Boh\'orquez Bar\'on, Diego Nicol\'as} \quad \texttt{dibohorquezb@unal.edu.co}\\[0.25em]
\textbf{M\'arquez Ariza, Sara Camila} \quad \texttt{samarqueza@unal.edu.co}\\[0.25em]
\textbf{Monta\~no Ria\~no, Samuel} \quad \texttt{samontanor@unal.edu.co}\\[0.25em]
\textbf{Mor\'e D\'iaz, Carlos Andr\'es} \quad \texttt{cmore@unal.edu.co}}
\end{center}
\vspace{1.2cm}

\section{Declaración de uso de inteligencia artificial}

Este proyecto se desarrolló \textbf{con ayuda de un modelo de inteligencia
artificial}, y esta bitácora documenta ese uso con trazabilidad completa, tal
como exige la política de la entrega. Conviene, sin embargo, precisar la
naturaleza de esa ayuda, porque define la autoría del trabajo:

\begin{itemize}
\item \textbf{La sustancia del proyecto provino del creador.} El software se
obtuvo mediante una infraestructura preparada y aportada por él: la estructura
de carpetas del proyecto, las especificaciones técnicas por módulo
(\texttt{INSTRUCCIONES\_CORE.md}, \texttt{INSTRUCCIONES\_THERMO.md},
\texttt{INSTRUCCIONES\_MECH.md}, \texttt{INSTRUCCIONES\_GUI.md}), los datos
termofísicos de referencia, las referencias bibliográficas y normativas
(Kern, 1950; estándares TEMA; el informe técnico de auditoría con los casos
benchmark), los casos de prueba de ejemplo y el material gráfico del programa.
\item \textbf{La IA se utilizó exclusivamente como modelo de programación:}
para escribir, organizar, depurar y empaquetar código, y para redactar
documentación técnica a partir de los materiales anteriores.
\item \textbf{Las decisiones fueron del creador.} Cada criterio de producto,
prioridad, corrección o aceptación de resultados fue decidido por él: qué
versión conservar, qué errores corregir, cómo presentar los dibujos, cuándo
un resultado se daba por válido. La IA proponía alternativas; el creador
elegía y validaba.
\end{itemize}

La información generada por la IA se verificó antes de incorporarse al código,
a los cálculos o al documento final (ver sección~4).

\section{Modelo de inteligencia utilizado}

Para evitar ambigüedades en la trazabilidad —el pliego menciona a modo de
ejemplo a Gemini— se aclara el modelo realmente empleado:

\begin{center}
\begin{tabular}{@{}p{4.2cm}p{9.6cm}@{}}
\toprule
\rowcolor{seccion}
\textbf{Aspecto} & \textbf{Descripción} \\
\midrule
Herramienta & \textbf{Hermes Agent}, agente conversacional de
\textbf{Nous Research}, con acceso a terminal, sistema de archivos y
control de versiones. \\
Modelo de lenguaje & \textbf{MiMo-V2.6-Pro} (Primalabs AI), el modelo que
opera el agente Hermes en las sesiones de este proyecto. \\
Modalidad de uso & Programación asistida: el creador describe el objetivo y
los criterios; la IA implementa, verifica con ejecución real y documenta. \\
Alcance & Generación de código, depuración, pruebas automáticas,
empaquetado del ejecutable, control de versiones y redacción técnica.
No sustituye el juicio de diseño ni la validación de ingeniería. \\
\bottomrule
\end{tabular}
\end{center}

\section{Infraestructura aportada por el creador}

El trabajo de la IA se apoyó en un ecosistema que el creador preparó,
documentó y sostuvo durante todo el proyecto:

\begin{enumerate}
\item \textbf{Carpeta del proyecto y guías de trabajo:} un repositorio con
instrucciones por fases (validaciones, core, cálculo, interfaz, mecánica),
que fijaban restricciones técnicas y de redacción en cada módulo.
\item \textbf{Datos:} la base termofísica (10 fluidos con rangos de validez,
5 materiales TEMA Clase C), los casos de ejemplo en formato \texttt{.bb} y
las tablas del informe técnico de auditoría.
\item \textbf{Referencias:} Kern (1950), los estándares TEMA en las ediciones
9.ª y 10.ª, y el artículo guía de Tan y Fok (2006).
\item \textbf{Material gráfico:} el icono de la aplicación y el criterio
visual de las vistas técnicas.
\item \textbf{Decisiones y validación:} revisiones comparativas entre
versiones, criterios de corrección y la verificación final de resultados.
\end{enumerate}

\section{Bitácora de interacciones (principales prompts)}

A continuación se registran las interacciones efectivamente realizadas
durante el desarrollo, organizadas por el orden natural del trabajo sobre el
código: \textbf{construcción}, \textbf{edición e integración} y
\textbf{verificación y corrección}. Los prompts se presentan en su forma
depurada ---fieles al contenido y la intención originales, redactados para su
lectura técnica--- junto con el resultado y la verificación aplicada.

\textbf{A. Fase de construcción del código.}
\begin{enumerate}
\item \textbf{Base de datos termofísica.} «Construir la base de datos del
programa: los fluidos de trabajo (densidad, viscosidad dinámica, calor
específico y conductividad térmica en función de la temperatura, con su
ventana de validez declarada) y los materiales (los cinco de la norma TEMA
Clase C, con conductividad, rugosidad y resistencia de ensuciamiento), todo
en unidades SI y con la fuente de cada dato documentada.»
\item \textbf{Estructura sizing vs rating.} «Separar estrictamente el
cálculo en dos módulos: \emph{sizing} (dimensionar un equipo nuevo,
iterando sobre $U$) y \emph{rating} (evaluar un equipo ya construido,
resolviendo las temperaturas de salida). Los datos de entrada y los
resultados deben quedar claramente diferenciados para que el usuario no
confunda cuándo se usa uno u otro.» \emph{(El creador fijó además que
la interfaz pregunte primero el modo de trabajo ---sizing, rating u otras---
y que los ejemplos se dividan en esas dos familias.)}
\item \textbf{Correlaciones del método de Kern.} «Implementar las
correlaciones de Kern para el cálculo térmico e hidráulico ---balance
de energía, LMTD con factor $F$ de Bowman, Nusselt en tubos y coraza,
coeficiente global por resistencias en serie y caídas de presión
(Ecs.~7.43, 7.44 y 7.45/7.46)--- junto con los criterios mecánicos de TEMA
Clase C, citando cada ecuación y su rango de validez.»
\item \textbf{Interfaz amigable.} «Diseñar una interfaz gráfica agradable
y fácil de usar: formulario separado por modo de trabajo (datos de sizing y
de rating en bloques distintos), límites de $\Delta p$ permisibles dentro
del sizing, un checklist de verificación del diseño adaptado al modo,
dibujos técnicos vinculados a la geometría calculada, mensajes de error
educativos y guardado de proyectos \texttt{.bb}.»
\item \textbf{Diseño mecánico TEMA y planos.} «Implementar el diseño
mecánico Clase C (carcasa, cabezales, deflectores, espesores y criterios de
boquilla) y las vistas paramétricas acotadas, vinculadas a la geometría
calculada, para que cada plano se redibuje solo con los datos del
proyecto.»
\item \textbf{Persistencia y ejemplos.} «Implementar guardar y abrir
proyectos con los campos obligatorios del formulario, y dejar los ejemplos
del programa en ese mismo formato para que puedan cargarse y recalcularse
sin edición.»
\item \textbf{Suite de pruebas automáticas.} «Cubrir con una suite la base
termofísica, el cálculo térmico e hidráulico, el diseño mecánico, el manejo
de errores educativo y la interfaz, de modo que ninguna edición del código
se acepte sin pasarla.»
\item \textbf{Optimización y entrega.} «Optimizar el programa (cierre
consistente del haz real, estabilidad numérica) y la documentación,
verificarlo con una suite de pruebas automáticas y empaquetarlo como
ejecutable autocontenido de Windows.»
\end{enumerate}

\textbf{B. Fase de edición e integración del código.}
\begin{enumerate}\setcounter{enumi}{8}
\item \textbf{Integración de versiones.} «Lo único que tienes que hacer es
poner la que yo tenía ---dibujos de todas las carcasas y cabezales, mejores
mensajes de advertencia--- y agregarle las novedades de la v2.1, no
reemplazarlo. Mucho ojo.» \emph{(Resultado: la versión del creador quedó
como base y las novedades se añadieron con parches quirúrgicos módulo por
módulo; ningún archivo propio se sustituyó en bloque.)}
\item \textbf{Propiedad intelectual.} «En todo lo que tenga código deja
muy discretamente que es de propiedad de Diego Nicolás Bohórquez Barón.»
\emph{(Resultado: marca discreta en cada archivo fuente,
\texttt{LICENSE.txt}, la ventana «Acerca de» y la firma de los
artefactos.)}
\item \textbf{Modos estrictos.} «Si entro a sizing que solo pueda hacer
sizing, y si entro a rating que solo pueda hacer rating.»
\emph{(Resultado: la interfaz habilita únicamente los campos del modo
activo y valida las acciones por modo.)}
\item \textbf{Estilo visual.} «Maneja contrastes de colores como morados,
carmesí y naranjas; algunas letras siguen siendo bastante pequeñas; elimina
los emojis, que parezca menos IA y sea moderno.» \emph{(Resultado: paleta
morado--carmesí--naranja, tipografía ampliada, sin emojis y acabado
moderno.)}
\item \textbf{Ayuda y mensajes educativos.} «En la pestaña Help sigue
saliendo lo del régimen laminar; actualízala para que refleje las
correlaciones actuales.» \emph{(Resultado: ayuda reescrita con Hausen,
Gnielinski y Kern--Graetz, bandas de incertidumbre por régimen y la
excepción del metanol presurizado.)}
\item \textbf{Control de deflectores.} «Agregar el control del
espaciamiento de deflectores, con la recomendación $0{,}3$--$1\,D_s$ y un
botón para el valor intermedio.» \emph{(Resultado: barra de espaciamiento y
botón $\frac{1}{2}D_s$ en el diseño mecánico, con validación del rango
TEMA.)}
\item \textbf{Organización de la entrega.} «Los anexos van en la carpeta
04, no en otras partes.» \emph{(Resultado: entrega estructurada en
\texttt{01\_codigo\_fuente}, \texttt{02\_ejecutable},
\texttt{03\_documento} y \texttt{04\_bitacora\_ia}, con cada artefacto en
su carpeta.)}
\item \textbf{Actualización del documento con el LaTeX del equipo.} «Con
este LaTeX reemplaza el documento final; puedes eliminar la carpeta de
imágenes viejas. Usa tectonic para revisar que no se corran imágenes y
tablas.» \emph{(Resultado: documento sustituido, imágenes nuevas
incorporadas y compilación verificada sin errores ni desbordes de figuras y
tablas.)}
\end{enumerate}

\textbf{C. Fase de verificación y corrección del código.}
\begin{enumerate}\setcounter{enumi}{16}
\item \textbf{Verificación con ejercicios de la literatura.} «Sustituir los
casos de servicio por ocho ejercicios resueltos de la literatura ---cuatro
de \emph{sizing} y cuatro de \emph{rating}---, compararlos en el documento
final en un apartado nuevo y dejar el caso de Tan y Fok en un apartado
aparte con el análisis de sus diferencias. Lo importante es demostrar que
el programa sirve: que no tiene errores de cálculo de Reynolds ni de
presión y que compara bien.» \emph{(Resultado: ocho casos de Janna, Kern,
Sinnott--Towler, Subramanian y Karimi, con tablas de comparación por
magnitud; documento final, secciones 6 y 7.)}
\item \textbf{Fuentes oficiales y temperatura de evaluación.} «No se pueden
usar fuentes como Wikipedia: solo fuentes oficiales, libros donde se muestren
las propiedades en función de la temperatura; además, revisa que las
correlaciones se estén evaluando a la temperatura que corresponde.»
\emph{(Resultado: base termofísica revalidada contra Incropera A.4/A.5,
Perry's/DIPPR, IAPWS y ASHRAE; auditoría de temperaturas: propiedades a la
temperatura media de cada lado, $T_w$ para el factor de pared y cada
correlación en su régimen.)}
\item \textbf{Justificación de las desviaciones.} «¿Pero justificaste los
errores? Ahora agrega una justificación para lo de Tan y Fok; necesito que
todo en el documento quede lo mejor explicado.» \emph{(Resultado: tabla
«Causa atribuida a cada desviación» para los ocho ejercicios, su equivalente
para Tan y Fok y la descomposición cuantitativa de su diferencia: ninguna
desviación queda sin causa documentada.)}
\item \textbf{Caídas de presión y fuentes inconsistentes.} «No sé por qué
siguen existiendo casos donde el error asociado con la caída de presión es
tan abrupto, caso Kharaji; revisa los cálculos o dime el porqué. Además, lo
de Tan y Fok sigue siendo bastante ilógico, no entiendo por qué no da.»
\emph{(Resultado: el artículo de Karimi no publica su diámetro interno; su
velocidad exige $d_i = 9{,}244$ mm y con él la desviación de $\Delta p_t$
baja del $85\%$ al $5{,}7\%$; las diferencias restantes son de convención y
de la propia inconsistencia de las fuentes, todo documentado.)}
\item \textbf{Optimización del ejecutable.} «¿Puedo confiar en todo lo que
construiste? Porfa, optimiza el \texttt{.exe}, está algo lento.»
\emph{(Resultado: PyInstaller incluía PIL sin usar; excluido, el arranque
pasó de 12--25 s a 2,5 s (\emph{onefile}) y 0,35 s (versión rápida); la
confianza se respaldó con la suite de 153 pruebas y los scripts de
verificación incluidos en la entrega.)}
\item \textbf{Cálculo de tubos necesarios.} «Sigo percibiendo errores en el
cálculo de tubos necesarios; revísalo muy bien, ¡en uno está casi al
doble!» \emph{(Resultado: el punto de diseño podía quedar por debajo de la
ventana de velocidad e inflar el área; el \emph{sizing} ahora elige el
número de pasos para que la velocidad entre en la ventana ---como diseña
Kern--- y el caso afectado pasa de 348 a 268 tubos, con $U$ en el valor
publicado.)}
\item \textbf{Auditoría del Reynolds.} «Revisa muy bien los cálculos del
Reynolds; he notado que a veces es muy distinto entre teoría y práctica.
¿Por qué difieren? ¿Se basa en las iteraciones?» \emph{(Resultado: el
cálculo es determinista, no depende de las iteraciones; la diferencia viene
de la temperatura de evaluación de la viscosidad ---desplazamientos de 0,5 a
3,5 K explican del 1 al 9\%--- y, en una fuente, de calcular el $Re$ con el
diámetro exterior en vez del interior.)}
\item \textbf{Correlación 9 y factor de fricción.} «La correlación 9 tiene
dos formas, la de calentamiento y la de enfriamiento; aparte, revisa los
cálculos de la caída de presión ---es fundamental que funcione a la
perfección--- y cómo se calcula el factor de fricción. Creo que la cita de
la sección 6.4 está mal.» \emph{(Resultado: Dittus--Boelter implementado con
sus dos exponentes, $n = 0{,}4$ calentamiento y $n = 0{,}3$ enfriamiento
(Incropera, ec.~6.61); factores de fricción verificados frente a
Colebrook--White (error máximo $1{,}34\%$) y a los ejemplos publicados de
Kern ($+0{,}1\%$ y $+2{,}1\%$); la cita de Hewitt corregida a 1994 con sus
páginas y una fórmula del documento reparada.)}
\item \textbf{Método de coraza Bell--Delaware.} «¿Bell-Delaware es para cambios
de fase? Podrías cambiar/implementar lo de rigor; impleméntalo también en el
código y mételo en el documento final, y revisa que los resultados de los
ejemplos sean correctos.» \emph{(Resultado: módulo \texttt{bell\_delaware.py}
con las cinco corrientes, los factores $J_c$, $J_l$, $J_b$ y $J_r$ y la caída
de presión en tres componentes; selección automática del método de coraza (Kern o Bell--Delaware), 
con el método registrado en cada proyecto \texttt{.bb}; la sección 5.16
del documento lo describe y la tabla comparativa frente a Kern se verificó en
los cinco equipos de ejemplo; 12 pruebas nuevas en la suite, 165 en total.)}
\end{enumerate}

\section{Verificaciones, modificaciones y correcciones realizadas}

El equipo (creador asistido por IA) aplicó las siguientes verificaciones,
todas con evidencia reproducible:

\begin{enumerate}
\item \textbf{Suite de verificación automática (\texttt{verificar.py}):}
153 pruebas que cubren la base termofísica, el cálculo térmico e hidráulico,
el diseño mecánico TEMA, el manejo de errores educativo, la interfaz
gráfica, los proyectos de ejemplo y un barrido de robustez de 216
combinaciones de entrada. Cada modificación del código se validó contra
esta suite antes de aceptarse (resultados: 123/123 $\rightarrow$ 135/135
$\rightarrow$ 153/153 tras la revisión de regímenes).
\item \textbf{Auditoría de los errores documentados:} los 9 conflictos del
informe técnico se contrastaron con el código fuente y se verificó su
corrección efectiva (o su corrección puntual cuando persistían), con
pruebas dedicadas para cada uno.
\item \textbf{Corrección de regresiones propias:} se detectó que un comando
de sincronización había incorporado archivos binarios generados al
repositorio; se retiraron, se agregó \texttt{.gitignore} y se cerró el
issue correspondiente con evidencia.
\item \textbf{Verificación real del ejecutable:} el \texttt{.exe} se
ejecutó en cada compilación, comprobando proceso vivo y cierre limpio; no
se dio por bueno un build por el simple hecho de terminar sin error.
\item \textbf{Trazabilidad por control de versiones:} cada cambio quedó
registró en commits con hash verificable (estado local = remoto).
\item \textbf{Verificación física:} ninguna correlación ni constante se
incorporó por invento; solo se usaron las ecuaciones y datos anclados en
Kern, TEMA y el informe técnico, con sus fuentes citadas.
\item \textbf{Cobertura completa de los regímenes laminar y de transición:}
los errores que aparecían cuando el cálculo caía en régimen laminar se
resolvieron con correlaciones clásicas verificadas en la documentación
técnica consultada en línea (guía de transferencia de calor de
neqsim/equinor): la solución de Graetz con entrada térmica de
\textbf{Hausen (1959)} para tubos en laminar, la de \textbf{Gnielinski
(1976)} para la transición ($2300 \le Re < 5\cdot 10^6$, con factor de
fricción de Petukhov) y la rama laminar de Kern con dependencia de Graetz
($D_e/B$) para la coraza. La continuidad de $Nu(Re)$ en los bordes de cada
banda se verificó numéricamente y un barrido de 504 combinaciones de
fluidos y caudales laminares no produjo ninguna excepción. Con estas
correlaciones el número de tubos del dimensionamiento baja en los casos de
transición ---el caso base pasa de $502$ a $470$ tubos--- y las desviaciones
frente al benchmark documentado del proyecto caen del $4$--$9\%$ al
$0{,}5$--$3{,}4\%$.
\item \textbf{Verificación de las caídas de presión frente a la literatura:}
las ecuaciones implementadas reproducen las ecuaciones 7.44--7.46 de Kern
(1950), incluida su convención de \emph{cuatro cabezas de velocidad por paso}
en la pérdida por retorno (cita literal del libro, verificada en el texto
completo consultado en línea). Como contraste numérico, aplicadas al equipo
de Tan y Fok con la velocidad que impone su artículo ($v = 0{,}6$ m/s) dan
$\Delta p_t = 1873$ Pa frente a sus $1975{,}59$ Pa publicados (desviación
$5{,}2\%$). Se cuantificó además la dispersión de la literatura en esa
convención (Serth $(2n_p-1)$ cabezas, Frank $2{,}5$, Bell--Mueller $3$,
HEXTRAN $2$, Saunders $0{,}9$--$1{,}6$ por paso): elegir otra convención
cambia $\Delta p_t$ entre un $11{,}5$ y un $37{,}8\%$, y esa diferencia de
criterio ---no un error de cálculo--- explica las discrepancias aparentes
frente a otras fuentes. La verificación reproducible queda en
\texttt{validar\_caidas.py} y la de correlaciones en
\texttt{validar\_literatura.py}.

\item \textbf{Auditoría final de correlaciones, temperaturas y fricción:}
se verificó que cada correlación se evalúa a su temperatura propia
(propiedades a la temperatura media de cada lado, $T_w$ para el factor de
pared), que la forma de Dittus--Boelter aplica sus dos exponentes según el
sentido del calor ($0{,}4/0{,}3$) y que los factores de fricción reproducen
a Colebrook--White (error máximo $1{,}34\%$ para la aproximación de Haaland)
y los ejemplos publicados de Kern ($+0{,}1\%$ en $\Delta p_t$ y $+2{,}1\%$
en $\Delta p_s$). El Reynolds se descompuso frente a la literatura
evaluando la viscosidad a distintas temperaturas, lo que descartó a las
iteraciones como causa de las diferencias.
\end{enumerate}

\section{Alcance y limitaciones declaradas}

\begin{itemize}
\item La IA es una herramienta de \textbf{programación y redacción}; no
certifica diseños ni sustituye el análisis de ingeniería. Los resultados de
la herramienta son educativos y no constituyen planos de fabricación.
\item Cualquier resultado numérico generado durante el desarrollo fue
ejecutado y comparado antes de darse por bueno; las cifras de este documento
y del informe salen de ejecuciones reales del programa.
\item Las decisiones de producto —qué conservar, qué corregir, cómo
presentar— fueron del creador en todas las interacciones registradas.
\end{itemize}

\vspace{0.6cm}
\noindent{\color{acento}\rule{\linewidth}{2pt}}\\[0.3cm]
{\small\color{gristexto}\raggedright Documento generado como anexo de trazabilidad del uso
de inteligencia artificial. Proyecto: Intercambiador de Tubos y Coraza
(método de Kern / TEMA Clase C).}

\section*{Anexo: Procedencia de la base de datos termof\'isicas}
\addcontentsline{toc}{section}{Anexo: Procedencia de la base de datos}

La infraestructura de datos ---fluidos, materiales y ensuciamientos---
es obra del creador del proyecto; la IA la documenta aqu\'i con sus
fuentes, tal como consta en el c\'odigo fuente.


Este anexo documenta el origen de cada dato de la base termof\'isica del
programa (\texttt{src/core/thermo\_data.py}), la \'unica fuente de
verdad del motor de c\'alculo. Todas las correlaciones se eval\'uan en
unidades SI con la temperatura en kelvin, declaran su ventana de validez y
\textbf{rechazan el c\'alculo fuera de ella} (error \texttt{ValueError})
en lugar de extrapolar: una correlaci\'on emp\'irica extrapolada no es
f\'isica.

\subsection*{C.1. Fluidos: ventanas de validez y fuentes del ajuste}

\begin{center}
{\small
\begin{tabular}{@{}lll@{}}
\toprule
\textbf{Fluido} & \textbf{Ventana declarada} & \textbf{Fuente del ajuste} \\
\midrule
Agua l\'iquida & $0$--$100\,^{\circ}$C & IAPWS-IF97 (l\'iquido saturado) \\
Vapor de agua ($1$ atm) & $100$--$600\,^{\circ}$C & Incropera, Tabla A.4 \\
Aire ($1$ atm) & $-73$--$727\,^{\circ}$C & Incropera Tabla A.4; $\mu$ y $k$ por Sutherland \\
Etanol & $0$--$75\,^{\circ}$C & Perry's, octava ed. / tipo DIPPR \\
Metanol & $0$--$60\,^{\circ}$C & \textit{idem} \\
Etilenglicol & $0$--$197\,^{\circ}$C & \textit{idem} \\
Benceno & $10$--$80\,^{\circ}$C & \textit{idem} \\
Tolueno & $-20$--$100\,^{\circ}$C & \textit{idem} \\
n-Hexano & $-20$--$60\,^{\circ}$C & \textit{idem} \\
Amon\'iaco l\'iquido & $-40$--$50\,^{\circ}$C & ASHRAE / Perry's (l\'iquido saturado) \\
\bottomrule
\end{tabular}}
\end{center}

\subsection*{C.2. Ajuste y validaci\'on}

Los polinomios de $\rho$, $c_p$ y $k$, y las correlaciones de viscosidad,
se ajustaron por m\'inimos cuadrados a las tablas de las fuentes citadas.
Dentro del rango declarado el error es menor del $2\%$ en $\rho$, $c_p$ y
$k$ (t\'ipicamente menor del $0\,5\%$); la viscosidad ---la propiedad m\'as
sensible a la fuente tabulada--- puede desviarse hasta un $5\%$ en
alcoholes y glicol y un $20\%$ en amon\'iaco l\'iquido. La validaci\'on
cuantificada se genera con \texttt{src/core/informe\_propiedades.py}
(\texttt{informe\_propiedades.pdf}) y la suite autom\'atica la reverifica
(p.~ej., agua a $25\,^{\circ}$C: $c_p = 4181$ J/kg$\cdot$K de referencia).

\subsection*{C.3. Materiales s\'olidos (TEMA Clase C)}

\begin{center}
{\small
\begin{tabular}{@{}lrrrr@{}}
\toprule
\textbf{Material} & \textbf{$k_m$ (W/m$\cdot$K)} & \textbf{$\rho_m$ (kg/m\textsuperscript{3})} & \textbf{$\varepsilon$ ($\mu$m)} & \textbf{$R_f$ (m\textsuperscript{2}K/W)} \\
\midrule
Cobre & 386,0 & 8933 & 1,5 & $1\,761\times 10^{-4}$ \\
Acero al carbono & 60,5 & 7854 & 46 & $3\,522\times 10^{-4}$ \\
Acero inoxidable AISI 304 & 14,9 & 7900 & 1,5 & $1\,761\times 10^{-4}$ \\
Lat\'on Admiralty & 111,0 & 8530 & 1,5 & $1\,761\times 10^{-4}$ \\
Aluminio & 237,0 & 2702 & 1,5 & $1\,761\times 10^{-4}$ \\
\bottomrule
\end{tabular}}
\end{center}

Conductividad y densidad a unos $300$ K; la rugosidad distingue tubo
estirado ($1{,}5~\mu\mathrm{m}$) de acero comercial ($46~\mu\mathrm{m}$).

\subsection*{C.4. Factores de ensuciamiento}

Los factores siguen la tabla TEMA RGP-T-2.4, con la conversi\'on
$1~\mathrm{h\cdot ft^2\cdot ^\circ F/Btu} = 0{,}17611~\mathrm{m^2 K/W}$;
de ah\'i los $R_f = 1{,}761\times 10^{-4}$ y $3{,}522\times 10^{-4}$
m$^2$K/W de la tabla anterior. La jerarqu\'ia que aplica el programa es:
$R_f$ cargado por el usuario $>$ tabla del fluido $>$ mitad del $R_f$ del
material; el caso de referencia de Tan y Fok reproduce exactamente los
factores del art\'iculo ($R_{f,i} = 1{,}76\times 10^{-4}$ y
$R_{f,o} = 3{,}52\times 10^{-4}$ m$^2$K/W).


\subsection*{C.5. Base documental de los dibujos t\'ecnicos}

Las cinco vistas param\'etricas del programa (corte longitudinal,
distribuci\'on de tubos, placa deflectora, cabezal y brida, y el cat\'alog@
TEMA) se basan en el \emph{Standards of the Tubular Exchanger Manufacturers
Association}, d\'ecima edici\'on (TEMA, 2019) ---en particular su
Figura~N-1.2 de tipos de carcasas y cabezales--- y en el documento
\emph{TEMA Standards FAQ} (2020). Esos dos documentos, junto con el
\texttt{informe\_propiedades.pdf} que respalda la base termof\'isica, se
incluyen literalmente en la carpeta
\texttt{04\_bitacora\_ia/referencias/} de esta entrega, de modo que cada
dibujo y cada dato puedan rastrearse hasta su fuente primaria.

\end{document}
